\subsection{Attori}
Gli \textit{attori}\textsubscript{G} individuati sono distinti in primari, che avviano l'interazione con il sistema, e secondari, ovvero sistemi esterni che forniscono o ricevono servizi dal sistema.

\begin{itemize}
    \item \textbf{Mittente}: attore primario. Rappresenta il cliente \textit{B2B}\textsubscript{G}, identificato tramite \textit{anagrafica}\textsubscript{G}, che invia richieste non strutturate attraverso i canali monitorati. Il sistema opera in modo autonomo: le fasi della \textit{pipeline}\textsubscript{G} sono quindi modellate come \textit{casi d'uso}\textsubscript{G} innescati dalla richiesta del mittente, che non interagisce direttamente con l'interfaccia (VE-1.1, VE-1.4, LG-1.1).
    \item \textbf{Utente in sola lettura}: attore primario. \textit{Utente}\textsubscript{G} dell'interfaccia amministrativa che consulta richieste, \textit{log}\textsubscript{G} e \textit{dashboard}\textsubscript{G} di monitoraggio, senza poter modificare la configurazione o i dati (LG-3.3, LG-4.2).
    \item \textbf{Gestore}: attore primario. Utente dell'interfaccia amministrativa che, oltre alle funzionalità dell'Utente in sola lettura, gestisce le regole di azione e notifica e le richieste non riuscite (LG-3.3).
    \item \textbf{Amministratore}: attore primario. Utente dell'interfaccia amministrativa che, oltre alle funzionalità del Gestore, configura canali, applicativi di destinazione, modello \textit{LLM}\textsubscript{G}, posta in uscita e \textit{soglia di confidenza}\textsubscript{G}, gestisce anagrafica, utenti e profili e la persistenza dei dati (LG-3.3).
    \item \textbf{Canale di ingresso}: attore secondario (sistema esterno). Sorgente da cui il sistema acquisisce le richieste: casella di posta elettronica, canale WhatsApp o, come funzionalità desiderabile, cartella del filesystem, locale o remota (VE-1.1, LG-1.1).
    \item \textbf{Motore di estrazione (OCR / Speech-to-Text)}: attore secondario (sistema esterno). Componente che converte in testo audio, immagini e documenti scansionati, ad esempio Whisper, Google Speech, Tesseract o EasyOCR (VE-1.2, LG-1.1).
    \item \textbf{LLM}: attore secondario (sistema esterno). Fornisce i servizi di raffinamento dell'estrazione, individuazione dell'intenzione, estrazione delle entità e strutturazione dei dati tramite interfaccia compatibile con le \textit{API}\textsubscript{G} di OpenAI; è un modello a pesi aperti eseguito in locale, ad esempio Qwen (VE-1.8, LG-2.1, LG-2.2).
    \item \textbf{Applicativo aziendale di destinazione}: attore secondario (sistema esterno). Portale B2B, \textit{ERP}\textsubscript{G}, \textit{CRM}\textsubscript{G} o sistema di ticketing che riceve i dati strutturati tramite \textit{API REST}\textsubscript{G}; in fase di sviluppo la \textit{proponente}\textsubscript{G} fornisce \textit{mock}\textsubscript{G} e API dedicate (VE-1.1, LG-2.3).
    \item \textbf{Servizio di posta in uscita}: attore secondario (sistema esterno). Servizio \textit{SMTP}\textsubscript{G} utilizzato per l'invio di notifiche e messaggi al mittente (Capitolato).
\end{itemize}